\section{Minimal roadmap and kernel-relative posture}

This manuscript is \emph{kernel-relative}. It does not re-derive the kernel from below. Appendix~A supplies the full technical kernel floor that distinguishes bare trace checking from inferential governance, excludes weaker verdict regimes, fixes kernel attachment at inferential governance, and records local kernel closure. Relative to that floor, the interface, quotient, and coherent-reasoning layers are treated here as forced consequences of the same kernel rather than as parallel foundations.

The dependency order is strict:
\begin{enumerate}[leftmargin=2.5em]
  \item full technical Appendix~A floor;
  \item forced downstream interface, quotient, and coherent-reasoning consequences;
  \item inferential-life bridge;
  \item code-role exhaustion;
  \item diagonal-route failure and rescue-route exhaustion;
  \item terminal same-domain non-instantiability theorem.
\end{enumerate}

This ordering matters because the coherent-reasoning layer defeats the hardest hostile attack. The attack says that a same-domain G\"odel threat might matter only at hidden intermediate proof paths or proof-theoretic strength levels without surfacing in standing classification or coherent trajectory behavior. Once coherence is fixed as pointwise standing-preservation along legal same-scope trajectories, once no-deferred-coherence is in force, once terminal standing certifies coherence, and once all coherence-relevant distinctions factor through tensor content rather than skin, that hostile option is gone.


Detailed formalization ledgers, long hostile-referee crosswalks, and the hyper-armored dependency reservoir are intentionally removed from the primary manuscript package and retained in the audit companion. The main paper therefore keeps only what is needed for the body path and the full technical appendices A--E.

The visible body path is therefore exact and finite. The manuscript now asks the reader to track only the following eight results in order: inferential-life exhaustion, closure-relevance transmission, code-predicate role exhaustion, the necessary package for a same-domain Godel threat, no same-domain diagonalization, rescue-route exhaustion, exhaustion of same-domain incompleteness-objection routes, and the terminal non-instantiability theorem. Everything else has either been placed in the full technical appendices A--E or relegated to the audit companion.